\subsection{Twisted higher Chow groups and the image of the reduced norm}\label{sec:paper-global-weight-one}
In this subsection, we combine the local--global principle of Section~\ref{sec:paper-Q-local-global} with a cyclic splitting field to show that $\SK_2(D)=0$ over global fields, and we recover the formula of Bak and Rehmann \cite[\S2, Theorem 3]{BakRehmann1984} for the image of the reduced norm (Theorem~\ref{thm:paper-global-norm}).
We also introduce twisted higher Chow groups and prove the vanishing in weight one that is used in Section~\ref{sec:global-spectrum} (Theorem~\ref{thm:paper-global-weight-one}).


\begin{lemma}\label{lem:paper-cyclic-real-split}
Let $F$ be a global field and $D$ a central division algebra over $F$.
There is a cyclic maximal subfield $L_0\subset D$ that splits completely at every real place of $F$ at which $D$ splits.
\end{lemma}
\begin{proof}
If $F$ is a global function field or $\deg D$ is odd, a cyclic maximal subfield of $D$ has the required property, since a Galois extension of odd degree splits completely at all real places.
Let $F$ be a number field and $\deg D$ even.
Starting from a cyclic maximal subfield, we correct its class field theory character by a quadratic character that is trivial at every place where $D$ does not split and cancels the signs at the real places where $D$ splits; weak approximation of square classes gives these local conditions.
The corrected character has order dividing $\deg D$ because the original and quadratic characters do.
At the places where $D$ does not split, the local degrees are unchanged and are multiples of the local indices, whose least common multiple is $\deg D$.
Hence the corrected cyclic extension $L_0/F$ has degree $\deg D$ and splits $D$, so it embeds in $D$ as a maximal subfield.
\end{proof}

\begin{theorem}\label{thm:paper-global-norm}
For a global field $F$ and any central division algebra $D$ over $F$, we have $\SK_2(D)=0$.
If $S_{\mathbb R}$ is the set of ramified real places of $D$, then the image of the reduced norm is
\[
 R_D=\ker\left(K_2(F)\xrightarrow{(s_v)_{v\in S_{\mathbb R}}}
                        (\Z/2)^{S_{\mathbb R}}\right).
\]
Here $s_v\{a,b\}=1$ if both $a$ and $b$ are negative at $v$, and is zero otherwise.
Consequently $C_2(D/F)\simeq(\Z/2)^{S_{\mathbb R}}$.
\end{theorem}

We use cyclic Hilbert~90 with the field $L_0$ of Lemma~\ref{lem:paper-cyclic-real-split} to prove injectivity of $\Nrd_2$ on $I_F(D)$.
Local generation and Theorem~\ref{thm:paper-Q-local-global} then give $K_2(D)=I_F(D)$, whose reduced norm image is $N_{L_0/F}K_2(L_0)$.

\begin{proof}
Let $L_0/F$ be as in Lemma~\ref{lem:paper-cyclic-real-split}.
We first show that $\Nrd_2$ is injective on $I_F(D)$.
Every real place of a finite splitting field $E/F$ lies over a real place of $F$ at which $D$ splits, and $L_0$ splits completely there.
Hence no real place of $E$ becomes complex in $EL_0$, so the field norm $K_2(EL_0)\to K_2(E)$ is surjective \cite[\S2, Theorem 2]{BakRehmann1984}.
Transitivity of transfer gives
\[
 I_F(D)=\theta_{L_0}K_2(L_0).
\]
Let $\sigma$ be a generator of $G=\Gal(L_0/F)$.
By Skolem--Noether, $\theta_{L_0}\bigl((1-\sigma)K_2(L_0)\bigr)=0$, so $\theta_{L_0}$ factors through $K_2(L_0)_G$ and we obtain
\[
 K_2(L_0)_G\xrightarrow{\ \theta_{L_0}\ }I_F(D)\xrightarrow{\ \Nrd_2\ }K_2(F),
 \qquad \Nrd_2\circ\theta_{L_0}=N_{L_0/F}.
\]
If $\operatorname{char}F\nmid\deg D$, cyclic Hilbert 90 \cite[Theorem 14.1]{MerkurjevSuslin1982} gives $\ker N_{L_0/F}=(1-\sigma)K_2(L_0)$, so the composite is injective on $K_2(L_0)_G$.
If $\operatorname{char}F=p$ divides $\deg D$, let $P$ be a Sylow $p$-subgroup of $G$ and $M=L_0^P$.
Since $[L_0:L_0^p]=[M:M^p]=p$, the groups $K_2(L_0)$ and $K_2(M)$ are uniquely $p$-divisible \cite[Theorem 8.1]{GeisserLevine}, and the identities $\res_{L_0/M}\circ N_{L_0/M}=\sum_{g\in P}g$ and $N_{L_0/M}\circ\res_{L_0/M}=|P|$ give $K_2(L_0)_P\simeq K_2(M)$.
Cyclic Hilbert 90 for the cyclic extension $M/F$ of degree prime to $p$ then gives the same injectivity.
Since the first map is surjective, $\Nrd_2$ is injective on $I_F(D)$.

Finally we show that $K_2(D)=I_F(D)$.
At every finite place $v$, the group $K_2(D_v)$ is generated by the images of splitting transfers: for a Galois splitting field $L/F_v$ of $D_v$, the norm $N_{(D_v)_L/D_v}$ is $\theta_L$ under the split Morita identification, so Lemma~\ref{lem:paper-local-even-coinvariants} gives
\[
 \theta_L\colon K_2(L)_{\Gal(L/F_v)}\xrightarrow{\ \sim\ }K_2(D_v).
\]
At ramified real places, the splitting transfer $\theta_{\mathbb C}$ surjects onto $K_2$ of the real quaternion algebra as well \cite[Proposition 6.6.1(2)]{KahnLevine2010}.
Hence the target of Theorem~\ref{thm:paper-Q-local-global} vanishes, so $K_2(D)=I_F(D)$, and the injectivity of $\Nrd_2$ on $I_F(D)$ gives $\SK_2(D)=0$.
The image of $\Nrd_2$ is $N_{L_0/F}K_2(L_0)$, which equals $R_D$ by \cite[\S2, Theorem 2]{BakRehmann1984}.
Weak approximation realizes any real sign vector on a symbol $\{a,-1\}$, so the displayed sign map is surjective.
\end{proof}

\begin{definition}\label{def:paper-twisted-chow}
Let $X/F$ be a smooth equidimensional scheme, and put
\[
 D_X=D\otimes_F\mathcal O_X,\qquad \Delta_F^m=\Spec F[t_0,\ldots,t_m]\Big/\Bigl(\sum_i t_i-1\Bigr).
\]
An integral closed subscheme $W$ of codimension $r$ in $X\times\Delta_F^m$ is an admissible support if it meets every face $X\times\{t_{i_1}=\cdots=t_{i_s}=0\}$ properly.
We define
\[
 z^r(X,m;D)=\bigoplus_{W\ \text{admissible}}
                   K_0(D_{\kappa(W)}).
\]
Write $\partial_i^*$ for pullback along the $i$th face and $\partial=\sum_{i=0}^m(-1)^i\partial_i^*$ for the differential.
The pullbacks are constructed by $K$-theory with support and satisfy the face and degeneracy identities, so $z^r(X,\bullet;D)$ is a chain complex.
We set
\[
 \CH^r(X,m;D)=H_m\bigl(z^r(X,\bullet;D)\bigr)
\]
and abbreviate this to $\CH^r(F,m;D)$ when $X=\Spec F$.
For $D=F$ these are Bloch's higher Chow groups \cite{LevineHigherChow}, \cite[Definitions 5.6.1 and 5.6.3]{KahnLevine2010}.

The reduced rank gives an injection
\[
 K_0(D_{\kappa(W)})\hookrightarrow\Z,
 \qquad\im=\ind(D_{\kappa(W)})\Z.
\]
It induces the \emph{reduced cycle map} $\rho$, which realizes the twisted cycle complex as a subcomplex of the ordinary cycle complex.
Under this identification the coefficient differential is face restriction with the usual intersection multiplicities.
For finite coefficients we use the derived tensor product of the complex, and we distinguish the homology of this derived tensor product from the quotient of integral homology by the coefficient integer \cite[Sections 5.6 and 5.9, Theorem 6.5.3]{KahnLevine2010}.
\end{definition}

The same methods give the following vanishing for twisted higher Chow groups in weight one: in degree two we use the local--global principle, and in degree three the field $L_0$ and the reduced cycle map.

\begin{theorem}\label{thm:paper-global-weight-one}
Let $F$ be a global field and let $D$ be a central division algebra over $F$.
Then $\CH^1(F,2;D)=\CH^1(F,3;D)=0$.
\end{theorem}

\begin{proof}
\emph{Step 1: identification with $Q_k(D)$.}
Over a perfect field $k$, we use the edge exact sequence of \cite[Theorem 6.7.1(2)]{KahnLevine2010},
\[
 0\to\CH^1(k,3;D)\to\CH^2(k,2;D)\xrightarrow{e_2}K_2(D)
 \to\CH^1(k,2;D)\to0.
\]
Transfers from splitting fields generate $\CH^2(k,2;D)$ \cite[Lemmas 6.6.2--6.6.4]{KahnLevine2010}: at each admissible closed point $x$, we choose a maximal subfield of the division algebra in the Brauer class of $D_{\kappa(x)}$, and pushforward from the resulting point gives the generator of the $K_0$ coefficient.
Over a splitting field $E$, the map $e_2$ is the isomorphism $\CH^2(E,2)\simeq K_2(E)$ \cite{NesterenkoSuslin}, compatible with transfer.
Thus
\[
 \im e_2=I_k(D),\qquad \CH^1(k,2;D)=Q_k(D),
\]
naturally in field extensions.

\emph{Step 2: imperfect fields.}
For a field $k$ of characteristic $p$, we first invert $p$ and pass to the perfect closure of $k$.
For a finite purely inseparable extension $M/k$, we have
\[
 \cor_{M/k}\circ\res_{M/k}=\res_{M/k}\circ\cor_{M/k}=[M:k].
\]
For cycles, the factor $[M:k]$ is the product of the support multiplicity and the residue degree.
For $K$-theory, we filter the local Artinian ring of the self-base-change by the powers of its radical and use additivity \cite{QuillenHigher}.
Twisted cycle complexes, $K$-theory, and the quotient by splitting transfers commute with filtered colimits of fields \cite{KahnLevine2010}.
Hence restriction to the perfect closure is an isomorphism after inverting $p$, and
\[
 \CH^1(k,2;D)[1/p]\simeq Q_k(D)[1/p].
\]

\emph{Step 3: the $p$-primary part.}
For the $p$-primary part, restriction and corestriction along an extension of degree prime to $p$ that splits the prime-to-$p$ component of the Brauer class reduce the problem to a $p$-primary class.
Since $[F:F^p]=p$ for local fields $F$ of characteristic $p$ and for global function fields $F$, the rank-$1$ Frobenius comparison of Proposition~\ref{prop:app-frobenius} makes the reduced cycle map in positive weights a quasi-isomorphism.
Hence the $p$-primary parts in weight one and degrees two and three vanish by the commutative calculation \cite[Proposition 6.8.1]{KahnLevine2010}.

\emph{Step 4: degree two over $F$.}
The proof of Theorem~\ref{thm:paper-global-norm} gives $Q_F(D)=0$.
In characteristic zero, Step~1 gives $\CH^1(F,2;D)\simeq Q_F(D)=0$; in positive characteristic, Step~2 gives the same after inverting $p$, and Step~3 shows that the $p$-primary torsion vanishes.

\emph{Step 5: degree three over $F$.}
The generation of the diagonal group $\CH^2(F,2;D)$ by transfers from splitting fields holds over any field, so this step takes place over $F$; only the edge exact sequence below requires a perfect field.
For splitting fields $E$, the isomorphism $\CH^2(E,2)\simeq K_2(E)$ is compatible with transfer \cite{NesterenkoSuslin}.
For the field $L_0$ of Lemma~\ref{lem:paper-cyclic-real-split}, the field norms $K_2(EL_0)\to K_2(E)$ for splitting fields $E$ are surjective, as in the proof of Theorem~\ref{thm:paper-global-norm}, so the pushforward
\[
 \psi\colon\CH^2(L_0,2)\longrightarrow\CH^2(F,2;D)
\]
is surjective.
Then $\psi$ factors through the coinvariants of $G=\Gal(L_0/F)$ because the structure map $\Spec L_0\to\Spec F$ is $G$-invariant, and $\rho\circ\psi=N_{L_0/F}$ is the compatibility of $\rho$ with pushforward \cite[Lemma 6.6.4]{KahnLevine2010}.
Hence cyclic Hilbert 90, combined in characteristic $p$ with the averaging argument in the proof of Theorem~\ref{thm:paper-global-norm}, makes $\rho\circ\psi$ injective on coinvariants, and
\[
 \rho\colon\CH^2(F,2;D)\to\CH^2(F,2)
\]
is injective.
When $F$ is perfect, under $\CH^2(L_0,2)\simeq K_2(L_0)$ and $\CH^2(F,2)\simeq K_2(F)$ we have
\[
 \Nrd_2\circ e_2\circ\psi=\Nrd_2\circ\theta_{L_0}=N_{L_0/F}=\rho\circ\psi,
\]
and the surjectivity of $\psi$ gives $\Nrd_2\circ e_2=\rho$ \cite[Proposition 6.6.5]{KahnLevine2010}.
Hence the injectivity of $\rho$ and the edge exact sequence give
\[
 \CH^1(F,3;D)=\ker e_2\subseteq\ker\rho=0.
\]
In positive characteristic, restriction to the perfect closure commutes with $\rho$ and is an isomorphism on $\CH^2(-,2;D)[1/p]$ and $\CH^2(-,2)[1/p]$, so the injectivity of $\rho$ over $F$ shown above gives its injectivity over the perfect closure after inverting $p$.
There the edge exact sequence and the identity $\Nrd_2\circ e_2=\rho$ give $\CH^1(F,3;D)[1/p]=0$, while the $p$-primary part vanishes by the comparison above.
\end{proof}
